\chapter{From balanced equations to matrix maps}
\label{learn:matrix_maps}
\begin{learningcard}{Bring forward and choose your next operation}
\textsf{\textbf{Question.}} How can two accounts become one operation?\par
\textbf{Required lessons.} \hyperref[learn:proof]{From convincing examples to proof}\par
\textbf{Helpful background.} Spatial arrows help interpret a vector; units determine whether a norm is a physical length.\par
\textbf{Learn to.} Compute matrix actions, compositions, and real dot products with compatible dimensions.\par
\textbf{Entry check.} Solve x+y=7 and 2x+y=11.\par
\hyperref[solution:entry:matrix_maps]{Check your answer}; repair the operation in \href{https://github.com/Oichkatzelesfrettschen/precalc_paper}{Companion: Precalculus}.
\end{learningcard}


\section{How can two accounts become one operation?}
The companion's balanced-equation lesson recovers two yields from two
accounts. Keep that question, and change the representation. Suppose two
bundles of one grade and one of another yield eleven measures, while one
of the first and two of the second yield ten. These are modern teaching
accounts. Each grade has a fixed yield in measures per bundle.
\lessoncue{Notice}{The same two yields occur in both accounts.}
Write their numerical values in an ordered column
\[
 v=\begin{pmatrix}x\\y\end{pmatrix},\qquad
 A=\begin{pmatrix}2&1\\1&2\end{pmatrix},\qquad
 Av=\begin{pmatrix}2x+y\\x+2y\end{pmatrix}.
\]
The entries of $A$ count bundles; multiplying bundle counts by yields
produces measures. An ordered column of real numbers is a coordinate
vector. Order carries meaning: exchanging $x$ and $y$ exchanges the grades.
A matrix is a rectangular array; here its two rows describe the accounts.

\section{Combine corresponding coordinates}
For columns of equal length, add corresponding entries and multiply every
entry by a scalar:
\[
 \begin{pmatrix}4\\3\end{pmatrix}+\begin{pmatrix}1\\-1\end{pmatrix}
 =\begin{pmatrix}5\\2\end{pmatrix},\qquad
 2\begin{pmatrix}4\\3\end{pmatrix}=\begin{pmatrix}8\\6\end{pmatrix}.
\]
Subtraction adds the negative column. The zero column leaves every entry
unchanged under addition. A linear combination such as $2v-w$ means scale
$v$, then subtract corresponding coordinates of $w$. In a physical model,
corresponding quantities must have compatible units before addition.

\section{Multiply a row by a column}
For a matrix with $m$ rows and $n$ columns, an input column must have $n$
entries. Multiply matching entries of a row and the input, then add; repeat
for each row. The output has $m$ entries. In symbols,
$(Av)_r=\sum_{s=1}^n A_{rs}v_s$: the summation sign asks you to add the
products as $s$ runs from one to $n$.
\lessoncue{Worked example}{Check the recovered yields in both accounts.}
Substitute $x=4,y=3$:
\[
 A\begin{pmatrix}4\\3\end{pmatrix}
 =\begin{pmatrix}2\cdot4+1\cdot3\\1\cdot4+2\cdot3\end{pmatrix}
 =\begin{pmatrix}11\\10\end{pmatrix}.
\]
Solving $Av=(11,10)^T$ asks which input gives that output. The superscript
$T$ transposes an array: rows become columns. Thus $(11,10)^T$ is a column.
For a rectangular example,
\[
 \begin{pmatrix}1&2&3\\4&5&6\end{pmatrix}^{T}
 =\begin{pmatrix}1&4\\2&5\\3&6\end{pmatrix}.
\]
Transpose changes positions, while applying a matrix computes new values.

\begin{figure}[htbp]
\centering
\begin{tikzpicture}[>=Stealth,font=\small]
\node[draw,align=center,text width=2.8cm] (input) at (0,0)
 {$v=(4,3)^T$\\yields\\measures per bundle};
\node[draw,align=center,text width=2.8cm] (accounts) at (4.5,0)
 {$Av=(11,10)^T$\\account amounts\\measures};
\node[draw,align=center,text width=2.3cm] (difference) at (8.8,0)
 {$BAv=1$\\difference\\measure};
\draw[->] (input)--node[above] {$A$} (accounts);
\draw[->] (accounts)--node[above] {$B$} (difference);
\draw[->] (input.south) -- ++(0,-0.7) -|
 node[pos=0.25,below] {$BA$: subtract rows, then multiply} (difference.south);
\end{tikzpicture}

\caption{Read each arrow as an operation. The first produces two account
amounts; the second takes their difference. The composite row computes
that difference directly from the yields. Units travel through the arrows.}
\end{figure}

\section{Compose operations in their application order}
Take the first account minus the second with the row $B=(1,-1)$.
Then $B(Av)=11-10=1$ measure. Distribute the subtraction before inserting
numbers: $(2x+y)-(x+2y)=x-y$. Therefore
\[
 BA=\begin{pmatrix}1&-1\end{pmatrix}
     \begin{pmatrix}2&1\\1&2\end{pmatrix}
   =\begin{pmatrix}1&-1\end{pmatrix},\qquad B(Av)=(BA)v.
\]
The rightmost matrix acts first. In general, if $A$ has size $m\times n$
and $B$ has size $p\times m$, then $BA$ has size $p\times n$.
Its entry in row $r$, column $s$ is the row $r$ of $B$ multiplied by
column $s$ of $A$. The inner sizes must agree. Here $AB$ is undefined:
$A$ expects two input coordinates, while $B$ produces one.

Even when both products exist, their order can matter. Let
\[
 C=\begin{pmatrix}1&1\\0&1\end{pmatrix},\quad
 D=\begin{pmatrix}2&0\\0&1\end{pmatrix}.
\]
Row-column multiplication gives $CD=\begin{pmatrix}2&1\\0&1\end{pmatrix}$
and $DC=\begin{pmatrix}2&2\\0&1\end{pmatrix}$.
Doubling the first coordinate before adding the second differs from doubling
after the addition. The identity matrix $I$, with ones on its diagonal and
zeros elsewhere, satisfies $Iv=v$.

\section{Measure length only after choosing a meaning}
For real columns of equal length define the Euclidean dot product and norm:
\[
 u\mathbin{\cdot}v=u^Tv=\sum_{s=1}^n u_sv_s,\qquad
 \|v\|=\sqrt{v^Tv}=\sqrt{\sum_{s=1}^n v_s^2}.
\]
For displacement coordinates on perpendicular axes in the same length unit,
the Pythagorean theorem identifies this norm as length. For unrelated
quantities, choose scales and a physical interpretation before calling a
coordinate norm a measured distance.
\lessoncue{Worked example}{Find a unit direction and a perpendicular one.}
For $v=(3,4)^T$, the squared norm is $9+16=25$, so $\|v\|=5$.
Divide each coordinate by five to obtain $q=(3/5,4/5)^T$ with norm one.
For $w=(-4,3)^T$, $v^Tw=-12+12=0$: the directions are perpendicular.
A zero vector has norm zero; division by its norm is undefined.

\section{Try and check}
\paragraph{1. Transfer the account rule.}
For $M=\begin{pmatrix}3&1\\1&2\end{pmatrix}$ and $v=(2,5)^T$,
find $Mv$ and check the difference of its accounts directly.
\paragraph{Solution.}
Multiply each row: $Mv=(3\cdot2+5,2+2\cdot5)^T=(11,12)^T$.
Subtract the outputs to get $-1$. Subtract the rows first to obtain
$(2,-1)$, then compute $2\cdot2-5=-1$. Both orders give the same difference.

\paragraph{2. Check dimensions and order.}
Using $C,D$ above and $u=(1,2)^T$, calculate $C(Du)$ and $D(Cu)$.
\paragraph{Solution.}
First $Du=(2,2)^T$, then $C(Du)=(4,2)^T$.
First $Cu=(3,2)^T$, then $D(Cu)=(6,2)^T$. The intermediate columns
have two entries, so both operations are defined, but their outputs differ.

\paragraph{3. Normalize and test perpendicularity.}
Find the norm of $u=(5,12)^T$ and a perpendicular unit vector.
\paragraph{Solution.}
The squared norm is $25+144=169$, giving norm thirteen.
The vector $w=(-12/13,5/13)^T$ has squared norm
$(144+25)/169=1$ and $u^Tw=(-60+60)/13=0$.

\section*{Ready for the next question?}
\label{readiness:matrix_maps}
Apply the matrix with rows (3,1) and (1,2) to (2,5), then subtract the second output from the first.
\par\hyperref[solution:ready:matrix_maps]{Read the worked readiness solution.} Continue when you can explain the operations and assumptions. Otherwise return to the method and its solved exercises.

\lessonreference
{Match row and column sizes; multiply corresponding entries and add.
Apply $A$ before $B$ in $BA$. For real coordinates, $\|v\|^2=v^Tv$.}
{Coordinate order and units carry meaning. Transpose swaps rows and columns.
Matrix multiplication can depend on order; normalization requires $v\ne0$.}
{Chapter~\ref{learn:linear_algebra} uses these operations to find independent
coordinates and eigenvectors. Chapter~\ref{learn:complex_numbers} extends
length calculations to complex coordinates.}

